\documentclass[12pt, letterpaper]{article}

\usepackage{amsmath}
\usepackage{booktabs}
\usepackage{amsthm}
\usepackage{graphicx}
\usepackage[margin=1in]{geometry}
\usepackage{hyperref}
\hypersetup{colorlinks = true, linkcolor = blue, citecolor=blue, urlcolor = blue}
\usepackage{natbib}
\usepackage{enumitem}
\usepackage{setspace}
\usepackage[]{lineno}
\linenumbers*[1]
\doublespace

\title{\bf Effects of Short-Term and Long-Term Fatigue on MLB Pitch Performance: A Pitch-Level Statcast Analysis}
\author{Yelin Kim\\[1ex]
Department of Statistics, University of Connecticut}
\date{\today}

\begin{document}
\maketitle

\begin{abstract}
The purpose of this study is to examine how short-term fatigue, long-term workload, and days of rest influence pitch outcomes in Major League Baseball (MLB). Using pitch-level Statcast data, I model four key performance metrics—expected weighted on-base average (xwOBA), whiff rate, home run probability, and hard-hit events—against three types of fatigue measures: previous-game pitch count, 5-game average pitch count, and 10-game average pitch count. Consistent with prior work such as Bradbury and Forman (2012), the results show that short-term fatigue has minimal influence on pitch performance, while long-term cumulative workload significantly worsens outcomes across all four models. Days of rest show little to no meaningful effect. These findings demonstrate that fatigue operates on different timescales, and that cumulative workload is a more important predictor of performance decline than the commonly emphasized rest-day intervals.
\end{abstract}

\section{Introduction}
% intro 내용

\section{Methods}
% data + model + variables

\section{Results}
% xwOBA, whiff, HR, hard-hit 모델 결과

\section{Discussion}
% 해석 + 기존 문헌 비교 + 기여점 + 함의

\section{Conclusion}
% 요약 정리

\bibliographystyle{asa}
\bibliography{citations}

\end{document}
